\documentclass{article}
\usepackage{hyperref}
\usepackage{amsmath}
\usepackage[a4paper]{geometry}
\usepackage{fancyhdr}
\pagestyle{fancy}
\lhead{Röntgenstrahlung}
\rhead{September 2025}
\begin{document}
\section{Röntgenstrahlung}
Bei der Entstehung der Röntgenstrahlung wird eine Glühkathode mit einer $U_A$ von $100\, \text{V}$ - $200\, \text{kV}$ betreieben. Die austretenden e$^-$ treffen daraufhin auf eine Anode, bestehend aus einem Material mit einer hohen Ordnungszahl. Beim Auftreffen werden die Elektronen rasant abgebremst, weshalb Röntgenstrahlung entsteht. Die Anode ist generell als Trapez geformt, so dass die Röntgenstrahlung in die freie Seite scheinen kann.
 
Dabei besteht die Röntgenstrahlung eigentlich aus zwei Strahlungen: der \emph{Bremsstrahlung} und der \emph{charakteristische Strahlung}
\begin{description} 
 \item[Bremsstrahlung] Es gibt ein kontinuierliches Spektrum für die Energie jedes $f < f_{max}$. Diese geben dementsprechend weniger Energie ab. Hier wird beim Abbremsen des Elektrons dessen $E_{kin}$ direkt in Photonen umgewandelt.
 \item[Charakteristische Strahlung] Die Röntgenstrahlung zeigt bei bestimmten Wellenlängen Peaks auf. Die Lage dieser ist für das Anodenmaterial charakteristisch.
 
Sie entstehen dadurch, dass die Elektronen der inneren Atomschalen des Anodenmaterials aus ihrer Bahn herausgeschlagen werden. Dadurch, dass ein Elektron aus einer höheren Schale nun in diese erstellte Lücke einspringt, wird viel Energie in Form von Röntgenstrahlung freigegeben.
 
Dies ist des Weiteren mit den \hyperref[Energieniveaus]{Energieniveaus} zu erklären. Diese zeigen dass das Atom nur bei den Energien, die die Photonen bei den Wellenlängen der charakteristischen Peaks hat, das Atom in ein geringeres Energienieau fallen kann.
 
Nun hat jedes Material andere Energieniveaus. 
\end{description} 
 
\subsection{Kurzwellige Grenze $\lambda_{\text{min}}$} 
Es gilt 
\[
 e \cdot U = h \cdot f_{max} 
\]
Es folgt, beim darüber nachdenken,
\[
 \lambda_{min} = \frac{h \cdot c}{e \cdot U} 
\] 
 
\end{document}